\section{One root contrast teaches the child feature needed by another}
\label{sec:parity-toy}

Take four equiprobable root classes encoded by the signs $(U,UV)$, where
$U,V$ are independent uniform signs. The hidden left child carries $U$ and the
right child carries $V$. A root can predict the first contrast from its left
child alone. Predicting the second requires a feature from each subtree.

An unambiguous two-level grammar realizes these classes with four symbols
per level and four equally likely productions per symbol. Label each symbol
by two bits. For a leaf pair $x,x'$, the bottom parent is $x\oplus x'$.
For hidden child symbols whose first bits are $u,v$, the root is
$(u,u\oplus v)$, with $U=(-1)^u$ and $V=(-1)^v$. At either level each
output has exactly four input pairs. Thus uniform sampling of the root
and of each production gives the same distribution as sampling four
independent uniform leaves and evaluating the tree from the bottom up.
In particular the two subtree observations are independent. This grammar
is a specific choice of rules in the RHM family.

The two bottom nodes use the same trainable scalar table $h$. Its sixteen
entries span all functions on one ordered leaf pair. For a generic subtree
let $H$ be its hidden first-bit sign. Use the inner product
$\langle h,g\rangle=\mathbb E[hg]$ over the sixteen equiprobable leaf
pairs and define
\[
 a=\langle h,H\rangle,\qquad R=\langle h,h\rangle,\qquad
 r=a/\sqrt R .
\]
The two root predictions are $b h_L$ for $sU$ and $w h_Lh_R$ for $UV$.
The easy head uses a fixed constant channel on the right. With $s=1$
both contrasts supply a learning signal. With $s=0$ only the interaction
does. Both settings use labels available at the root. We train by
population squared loss with unit Euclidean gradient-flow rates for
$b,w$ and the coordinates of $h$ in a basis orthonormal under the stated
inner product. Unit-rate gradient flow in the raw table entries would
multiply the bottom learning rate by $1/16$ in these coordinates.
Initialize $R(0)=1$, $r(0)=\delta\in(0,1)$, and $b(0)=w(0)=0$.
Sending $\delta$ to zero removes initial alignment while keeping the
child table at unit norm.

\paragraph{A prerequisite measured by freezing the child table.}
Independence gives $\mathbb E[UVh_Lh_R]=a^2$ and
$\mathbb E[h_L^2h_R^2]=R^2$. Let $D$ denote the fraction of the hard
contrast's unit variance explained by its actual prediction. Completing
the square gives
\begin{equation}
 D=1-\mathbb E[(UV-wh_Lh_R)^2]
  =r^4-(wR-r^2)^2,\qquad \max_w D=r^4 .
 \label{eq:toy-risk}
\end{equation}
Thus reaching $D=p$ requires $r\ge p^{1/4}$ along the positive-alignment
trajectories considered here, even if the root has unlimited time to
refit. For example, a frozen child at $r=1/2$ allows a root to explain
at most $1/16$ of the interaction variance. To explain half, even the
best root needs $r\ge 2^{-1/4}$.

\paragraph{Root-only gradient flow.}
Including the easy head, the total loss is
\begin{equation}
 2\mathcal L=s^2+1+b^2R+w^2R^2-2sba-2wa^2 .
 \label{eq:toy-objective}
\end{equation}
Using $\nabla_h a=H$ and $\nabla_h R=2h$ gives
\begin{align}
 \dot b&=sa-bR,& \dot w&=a^2-wR^2,&
 \dot h&=(sb+2wa)H-(b^2+2w^2R)h .
 \label{eq:toy-parameter-flow}
\end{align}
Taking twice the inner product of the last equation with $h$ gives
$\dot R=2b\dot b+4w\dot w$, hence $R=1+b^2+2w^2$.
To obtain the alignment dynamics, take the inner product of $\dot h$
with $(H-rh/\sqrt R)/\sqrt R$. Its radial term cancels, leaving
\begin{align}
 \dot r&=(1-r^2)(sb/\sqrt R+2wr),&
 \dot b&=R(sr/\sqrt R-b),&
 \dot w&=R^2(r^2/R-w).
 \label{eq:toy-reduced}
\end{align}
The easy head responds to a single power of $r$, and its feedback toward
the hidden sign survives even at $r=0$ if $sb\ne0$. The interaction
head responds to $r^2$ and its own feedback contains another power of $r$.

\begin{theorem}
\label{prop:toy-gap}
Let $T_c(q)$ be the first time $r=q$ and $T_p(p)$ the first time $D=p$,
for fixed $0<q,p<1$ and sufficiently small $\delta>0$. Both times are
finite. In each arm both have the following expansion as $\delta\downarrow0$,
with constants in the remainder allowed to depend on the threshold.
\[
\begin{array}{ll}
s=0: & \displaystyle
 T=\frac{1}{4\delta^2}+\frac52\log\frac{1}{\delta}+O(1),\\[5pt]
s=1: & \displaystyle
 T=\frac{1}{\lambda_+}\log\frac{1}{\delta}+O(1),
 \qquad \lambda_+=(\sqrt5-1)/2 .
\end{array}
\]
In either arm, if $q<p^{1/4}$ then
\[
 0<\frac{p^{1/4}-q}{3}\ \le\ T_p(p)-T_c(q)\ =\ O_{p,q}(1).
\]
In particular this applies to $q=p=1/2$.
\end{theorem}

\begin{proof}
\emph{Uniform control after the initial escape.}
Equation~\eqref{eq:toy-reduced} preserves $0\le r,b,w\le1$ and $r$ is
nondecreasing there. Indeed $R\ge1$ and the targets of the two damped
readout equations lie in $[0,1]$. In turn $1\le R\le4$ by the invariant,
so the solution exists for all time. Once $r\ge\varepsilon>0$,
$w$ is bounded below after one further time unit by
$(1-e^{-1})\varepsilon^2/4$, by comparison in the damped $w$ equation.
Thus $\dot r\ge 2wr(1-r^2)$ carries all such trajectories toward one.
The hard prediction also converges to its target. The bounded gradient
flow has only stationary accumulation points. The loss derivative is
minus the squared gradient norm. That squared norm is uniformly continuous on
the bounded trajectory, so its finite time integral forces it to zero. At any accumulation point $r=1$, and
$\dot b=\dot w=0$ gives $b\sqrt R=s$ and $wR=1$. Hence $D\to1$.

The time needed for any fixed $r<1$ or $D<1$ threshold after
$r\ge\varepsilon$ is uniformly bounded in the seed. To see the uniformity,
start the reduced flow at any point of the compact set
$[\varepsilon,1]\times[0,1]^2$ in $(r,b,w)$. Each such state is realized by
a table with $R=1+b^2+2w^2$, so the convergence argument applies from
every point. Each trajectory strictly exceeds the desired level in
finite time. Continuous dependence gives a neighborhood that exceeds it
at the same time, and a finite subcover bounds all these times.
For the root threshold, choose $\varepsilon<p^{1/4}$ so it cannot occur
before the escape. Also $\dot r\le3$, by~\eqref{eq:toy-reduced}.
Travel from $q$ to the necessary $p^{1/4}$ therefore takes at least
$(p^{1/4}-q)/3$. The uniform bound gives the other direction.

\emph{Both root contrasts.}
For $s=1$ we show that escape to any sufficiently small fixed $r$ has
logarithmic time. Monotonicity of $r$ and variation of constants in
the two readout equations imply $b(t)\le r(t)$ and $w(t)\le r(t)^2$,
since their past forcing targets $r/\sqrt R$ and $r^2/R$ in~\eqref{eq:toy-reduced}
are at most $r(t)$ and $r(t)^2$, respectively. Thus $R=1+O(r^2)$ near the
origin, and
\[
 \dot r=b+O(r^3),\qquad \dot b=r-b+O(r^3).
\]
Writing $x=r+\lambda_+ b$ gives
$\dot x=\lambda_+x+O(x^3)$ and $x(0)=\delta$ because
$\lambda_+^2+\lambda_+=1$. Moreover $r\le x\le(1+\lambda_+)r$.
Until $r$ reaches a sufficiently small fixed $\varepsilon$,
$\dot x\ge\lambda_+x/2$, so
$\int x^2\,dt\le (2/\lambda_+)\int x\,dx=O(1)$.
Integrating $d\log x/dt=\lambda_++O(x^2)$ gives the stated escape time.
The bounds after $\varepsilon$ transfer it to both learning times.

\emph{Hard contrast alone.}
With $s=0$, $b$ stays zero. Put $u=r^2$, so $u(0)=u_0=\delta^2$ and
\[
 R=1+2w^2,\qquad \dot w=R(u-Rw),\qquad \dot u=4wu(1-u).
\]
Monotonicity and the damped root equation again give $0\le w\le u$.
Fix small $u_1>0$ below both $q^2$ and $p^{1/2}$, and let $t_1$ be the
first time $u=u_1$. The first two terms of the slow trajectory are
$g(u)=u-4u^2$. To control the error, write $e=w-g(u)$ and set
\[
 B(w,u)=R(u-Rw)-4g'(u)wu(1-u).
\]
Direct substitution gives $\dot e=B(w,u)$,
$B(g(u),u)=O(u^3)$, and $\partial_wB=-1+O(u)$ uniformly for $0\le w\le u$.
Taking $u_1$ small, the mean value theorem and variation of constants
give $|e(t)|\le C u_0e^{-t/2}+Cu(t)^3$, using that $u$ increases.
The exponential term accounts for $w(0)=0$. For
$F(u)=-1/(4u)+(5/4)\log u$ this gives
\[
 \frac{d}{dt}F(u(t))=(1+5u)(1-u)\frac{w}{u}
  =1+O(u^2)+O(e^{-t/2}).
\]
The time integral of the error is bounded uniformly in $\delta$. Indeed
$k=w/u\in[0,1]$ obeys
$\dot k=R(1-Rk)-4k^2u(1-u)\ge1-C_1 k$ on this interval.
Thus $k$ is bounded below by a positive constant for $1\le t\le t_1$,
and $\dot u=4ku^2(1-u)$ controls $\int_1^{t_1}u^2\,dt$.
Integration gives $t_1=F(u_1)-F(u_0)+O(1)$. Substituting
$u_0=\delta^2$ gives the stated expansion. The bounded time after $u_1$
transfers this expansion to the child and hard-contrast thresholds.
\end{proof}

The comparison changes just the easy target. The hard observable, TTN,
gradient-flow metric, and initial child representation are the same in
both arms. The bounded within-run lag and the divergent between-arm
speedup answer different questions. The first measures the time from a specified
child alignment to a specified fit by the trained root. The
second measures how much sooner the same interaction is learned when an
additional root contrast trains its prerequisite. Both statements concern
population squared error. Sign accuracy is insensitive to the magnitude
of a correctly signed score.
